\documentclass[11pt]{article}

\usepackage[margin=0.85in]{geometry}
\usepackage{amsmath,amssymb,booktabs,longtable,array}
\usepackage{enumitem}
\usepackage{xcolor}
\usepackage{listings}
\usepackage{fancyhdr}
\usepackage{hyperref}

\definecolor{navy}{RGB}{24,55,95}
\definecolor{green}{RGB}{25,100,70}
\definecolor{codebg}{RGB}{246,248,250}
\definecolor{codecomment}{RGB}{90,110,95}
\hypersetup{colorlinks=true,linkcolor=navy,urlcolor=navy}

\pagestyle{fancy}
\fancyhf{}
\lhead{STAT452 Chapter 1}
\rhead{R Exercises and Notes}
\cfoot{\thepage}
\setlength{\headheight}{14pt}

\lstdefinelanguage{Rcustom}{
  language=R,
  morekeywords={lm,coef,fitted,resid,model.matrix,confint,predict,anova,
    df.residual,rstandard,cooks.distance,lower.tail,newdata,interval},
  sensitive=true
}
\lstset{
  language=Rcustom,
  basicstyle=\ttfamily\small,
  keywordstyle=\color{navy}\bfseries,
  commentstyle=\color{codecomment}\itshape,
  stringstyle=\color{green},
  backgroundcolor=\color{codebg},
  frame=single,
  rulecolor=\color{black!15},
  breaklines=true,
  columns=fullflexible,
  keepspaces=true,
  showstringspaces=false,
  xleftmargin=0.4em,
  xrightmargin=0.4em,
  aboveskip=0.6em,
  belowskip=0.7em
}

\newenvironment{exercise}[1]
 {\par\medskip\noindent\color{navy}\textbf{Exercise #1.}\color{black}\ }
 {\par\medskip}
\newenvironment{answer}[1]
 {\par\medskip\noindent\color{green}\textbf{Solution #1.}\color{black}\ }
 {\hfill$\square$\par\medskip}

\newcommand{\E}{\operatorname{E}}
\newcommand{\Var}{\operatorname{Var}}

\title{Chapter 1: Multiple Linear Regression in R\\
\large Syntax Notes, Mathematical Connections, and Ten Guided Exercises}
\author{STAT452 -- Applied Statistics for Engineers and Scientists II}
\date{}

\begin{document}

\maketitle
\tableofcontents
\bigskip

\section{What this worksheet covers}

This worksheet follows Chapter 1 of the course: multiple linear regression,
least squares, fitted values, residuals, coefficient inference, $R^2$, nested
model comparison, and residual diagnostics. All examples use \texttt{mtcars},
a data frame included with R, so no package or download is required.

We will mainly use these variables:
\begin{center}
\begin{tabular}{lll}
\toprule
R name & Meaning & Unit\\
\midrule
\texttt{mpg} & fuel efficiency & miles per US gallon\\
\texttt{wt} & vehicle weight & 1000 pounds\\
\texttt{hp} & gross horsepower & horsepower\\
\bottomrule
\end{tabular}
\end{center}

The main model will be
\[
  Y_i=\beta_0+\beta_1x_{i,\mathrm{wt}}
      +\beta_2x_{i,\mathrm{hp}}+\varepsilon_i,
\]
or, in matrix form,
\[
  \mathbf Y=\mathbf X\boldsymbol\beta+\boldsymbol\varepsilon.
\]

\section{R syntax you should know}

\subsection{The basic mental model}

R stores information in \emph{objects}. You create an object with
\texttt{<-}, then give that object to functions that calculate a result.

\begin{lstlisting}
x <- c(2, 4, 6, 8)   # Store four values in x
mean(x)              # Give x to the mean() function
sd(x)                # Compute its sample standard deviation
\end{lstlisting}

The symbol \texttt{\#} starts a comment. R ignores everything after it on the
same line.

\subsection{Common syntax patterns}

\begin{longtable}{>{\ttfamily}p{0.29\textwidth}p{0.63\textwidth}}
\toprule
\normalfont Syntax & Meaning\\
\midrule
\endhead
x <- 5 & Assign the value 5 to the name \texttt{x}.\\
c(2, 4, 6) & Combine values into a vector.\\
cars\$wt & Select the \texttt{wt} column from the data frame \texttt{cars}.\\
cars[1:5, ] & Select rows 1 through 5 and every column.\\
cars[, c("mpg", "wt")] & Select every row and two named columns.\\
mean(x) & Call a function; \texttt{x} is its argument.\\
sum(x\^2) & Square every element of \texttt{x}, then add the results.\\
sqrt(x) & Take square roots element by element.\\
length(x) & Count the elements of a vector.\\
nrow(cars) & Count observations in a data frame.\\
head(cars) & Display the first six rows.\\
str(cars) & Display the structure and variable types.\\
summary(cars) & Produce quick summaries of every variable.\\
\bottomrule
\end{longtable}

\subsection{Regression formula syntax}

The R formula
\begin{lstlisting}
mpg ~ wt + hp
\end{lstlisting}
means ``model \texttt{mpg} using an intercept, \texttt{wt}, and \texttt{hp}.''
It corresponds to
\[
  \mathrm{mpg}_i=\beta_0+\beta_1\mathrm{wt}_i
                  +\beta_2\mathrm{hp}_i+\varepsilon_i.
\]

\begin{center}
\begin{tabular}{ll}
\toprule
R formula & Meaning\\
\midrule
\texttt{y \textasciitilde{} x} & intercept plus one predictor\\
\texttt{y \textasciitilde{} x1 + x2} & intercept plus two predictors\\
\texttt{y \textasciitilde{} .} & use every other column as a predictor\\
\texttt{y \textasciitilde{} x - 1} & omit the intercept\\
\texttt{y \textasciitilde{} x1:x2} & interaction only\\
\texttt{y \textasciitilde{} x1*x2} & both main effects and their interaction\\
\texttt{y \textasciitilde{} x + I(x\^2)} & linear and quadratic terms\\
\bottomrule
\end{tabular}
\end{center}

\noindent\textbf{Important:} In a formula, \texttt{+} means ``include another
term.'' It does not tell R to numerically add the two columns together.

\section{Translation between mathematics and R}

\begin{longtable}{p{0.30\textwidth}p{0.27\textwidth}p{0.35\textwidth}}
\toprule
Statistical object & Mathematics & R\\
\midrule
\endhead
Fit the model & minimize $\sum e_i^2$ & \texttt{fit <- lm(y \textasciitilde{} x1 + x2, data=d)}\\
Coefficient estimates & $\widehat{\boldsymbol\beta}$ & \texttt{coef(fit)}\\
Design matrix & $\mathbf X$ & \texttt{model.matrix(fit)}\\
Fitted values & $\widehat{\mathbf Y}=\mathbf X\widehat{\boldsymbol\beta}$ & \texttt{fitted(fit)}\\
Residuals & $\mathbf e=\mathbf Y-\widehat{\mathbf Y}$ & \texttt{resid(fit)}\\
Residual sum of squares & $\mathrm{SSE}=\sum e_i^2$ & \texttt{sum(resid(fit)\^2)}\\
Error degrees of freedom & $n-p-1$ & \texttt{df.residual(fit)}\\
MSE & $\mathrm{SSE}/(n-p-1)$ & \texttt{deviance(fit)/df.residual(fit)}\\
Residual standard error & $\sqrt{\mathrm{MSE}}$ & \texttt{summary(fit)\$sigma}\\
$R^2$ & $1-\mathrm{SSE}/\mathrm{SST}$ & \texttt{summary(fit)\$r.squared}\\
Coefficient table & estimate, SE, $t$, $p$ & \texttt{coef(summary(fit))}\\
Confidence intervals & $\widehat\beta_j\pm t^*\mathrm{SE}$ & \texttt{confint(fit)}\\
Prediction & $\widehat y(\mathbf x_0)$ & \texttt{predict(fit, newdata=...)}\\
Nested-model test & partial $F$-test & \texttt{anova(reduced, full)}\\
Diagnostic plots & assumption checks & \texttt{plot(fit)}\\
\bottomrule
\end{longtable}

\subsection{Matrix operations in R}

The following commands translate the matrix least-squares formula directly:
\begin{center}
\begin{tabular}{lll}
\toprule
Mathematics & R & Purpose\\
\midrule
$\mathbf X^{\mathsf T}$ & \texttt{t(X)} & transpose\\
$\mathbf X^{\mathsf T}\mathbf X$ & \texttt{t(X) \%*\% X} & matrix multiplication\\
$(\mathbf X^{\mathsf T}\mathbf X)^{-1}$ & \texttt{solve(t(X) \%*\% X)} & inverse\\
$\widehat{\boldsymbol\beta}$ & \texttt{solve(t(X) \%*\% X, t(X) \%*\% Y)} & solve normal equations\\
\bottomrule
\end{tabular}
\end{center}

Use \texttt{\%*\%} for matrix multiplication. Ordinary \texttt{*} performs
element-by-element multiplication. For numerical work, \texttt{solve(A, b)}
is preferable to \texttt{solve(A) \%*\% b} because it solves
$\mathbf A\mathbf x=\mathbf b$ without explicitly forming the inverse.

\subsection{How to read \texttt{summary(lm(...))}}

The coefficient table contains four columns:
\begin{enumerate}[leftmargin=*]
 \item \textbf{Estimate:} the fitted coefficient $\widehat\beta_j$.
 \item \textbf{Std. Error:} estimated standard deviation of
 $\widehat\beta_j$ across repeated samples.
 \item \textbf{t value:} $t=\widehat\beta_j/\mathrm{SE}(\widehat\beta_j)$
 for testing $H_0:\beta_j=0$.
 \item \textbf{Pr($>|t|$):} the two-sided $p$-value.
\end{enumerate}

The last few lines give the residual standard error, $R^2$, adjusted $R^2$,
and overall $F$-test. A small coefficient $p$-value concerns one predictor
conditional on the others; the overall $F$-test concerns all slopes together.

\newpage
\section{Exercises}

Start every exercise with:
\begin{lstlisting}
data(mtcars)
cars <- mtcars
\end{lstlisting}

\begin{exercise}{1: Inspect the data}
Use \texttt{head()}, \texttt{dim()}, \texttt{names()}, and \texttt{str()} to
inspect \texttt{cars}.
\begin{enumerate}[label=(\alph*),leftmargin=*]
 \item How many observations and variables are present?
 \item Are \texttt{mpg}, \texttt{wt}, and \texttt{hp} stored as numerical
 variables?
 \item What does one row represent?
\end{enumerate}
\end{exercise}

\begin{exercise}{2: Summaries and the simple-regression formula}
Let $X=\texttt{wt}$ and $Y=\texttt{mpg}$.
\begin{enumerate}[label=(\alph*),leftmargin=*]
 \item Calculate $\bar x$, $\bar y$, $s_x$, and $s_y$ with R.
 \item Calculate
 \[
  S_{xx}=\sum(x_i-\bar x)^2,
  \qquad
  S_{xy}=\sum(x_i-\bar x)(y_i-\bar y).
 \]
 \item Use $\widehat\beta_1=S_{xy}/S_{xx}$ and
 $\widehat\beta_0=\bar y-\widehat\beta_1\bar x$.
\end{enumerate}
\emph{Syntax hint:} multiplication is \texttt{*}, powers use \texttt{\^},
and \texttt{sum()} adds a vector.
\end{exercise}

\begin{exercise}{3: Simple linear regression}
Fit \texttt{mpg \textasciitilde{} wt} using \texttt{lm()}.
\begin{enumerate}[label=(\alph*),leftmargin=*]
 \item Extract the coefficients with \texttt{coef()}.
 \item Write the fitted equation.
 \item Interpret the slope, remembering that \texttt{wt} is measured in
 thousands of pounds.
 \item Confirm that the coefficients agree with Exercise 2.
\end{enumerate}
\end{exercise}

\begin{exercise}{4: Multiple linear regression}
Fit \texttt{mpg \textasciitilde{} wt + hp}.
\begin{enumerate}[label=(\alph*),leftmargin=*]
 \item Write the fitted equation.
 \item Interpret the \texttt{wt} coefficient while holding horsepower fixed.
 \item Interpret the \texttt{hp} coefficient while holding weight fixed.
 \item Compare the \texttt{wt} coefficient with the simple-regression value.
 Why might it change?
\end{enumerate}
\end{exercise}

\begin{exercise}{5: Fitted values and residual identities}
For the model in Exercise 4, store
\begin{lstlisting}
y_hat <- fitted(full_model)
e <- resid(full_model)
\end{lstlisting}
Then verify numerically that
\[
 y_i=\widehat y_i+e_i,
 \qquad
 \sum e_i=0,
 \qquad
 \sum \mathrm{wt}_i e_i=0,
 \qquad
 \sum \mathrm{hp}_i e_i=0.
\]
Explain why R may display numbers such as $10^{-13}$ instead of exact zero.
\end{exercise}

\begin{exercise}{6: Matrix least squares}
Use
\begin{lstlisting}
X <- model.matrix(full_model)
Y <- matrix(cars$mpg, ncol = 1)
\end{lstlisting}
and calculate
\[
 \widehat{\boldsymbol\beta}
 =(\mathbf X^{\mathsf T}\mathbf X)^{-1}
   \mathbf X^{\mathsf T}\mathbf Y.
\]
Compare the result with \texttt{coef(full\_model)}. Also calculate
\texttt{t(X) \%*\% e}. What does the result mean?
\end{exercise}

\begin{exercise}{7: SSE, MSE, and $R^2$}
For the two-predictor model, manually calculate
\[
 \mathrm{SSE}=\sum e_i^2,
 \qquad
 \mathrm{MSE}=\frac{\mathrm{SSE}}{n-p-1},
 \qquad
 R^2=1-\frac{\mathrm{SSE}}{\mathrm{SST}}.
\]
Use $n=32$ and $p=2$. Compare your results with
\texttt{summary(full\_model)}.
\end{exercise}

\begin{exercise}{8: A coefficient test and confidence interval}
For the \texttt{hp} coefficient, test
\[
 H_0:\beta_{\mathrm{hp}}=0
 \qquad\text{against}\qquad
 H_1:\beta_{\mathrm{hp}}\ne0.
\]
\begin{enumerate}[label=(\alph*),leftmargin=*]
 \item Extract its estimate and standard error.
 \item Calculate the $t$ statistic manually.
 \item Calculate the two-sided $p$-value using \texttt{pt()}.
 \item Obtain a 95\% confidence interval using \texttt{confint()}.
 \item State the conclusion at $\alpha=0.05$.
\end{enumerate}
\end{exercise}

\begin{exercise}{9: Confidence versus prediction intervals}
For a car with \texttt{wt = 3} and \texttt{hp = 110}, calculate:
\begin{enumerate}[label=(\alph*),leftmargin=*]
 \item the predicted fuel efficiency;
 \item a 95\% confidence interval for the mean fuel efficiency;
 \item a 95\% prediction interval for one new car.
\end{enumerate}
Explain why the prediction interval is wider.
\end{exercise}

\begin{exercise}{10: Nested models and diagnostics}
Let the reduced model be \texttt{mpg \textasciitilde{} wt} and the full model
be \texttt{mpg \textasciitilde{} wt + hp}.
\begin{enumerate}[label=(\alph*),leftmargin=*]
 \item Use \texttt{anova(reduced\_model, full\_model)} to test whether adding
 \texttt{hp} significantly improves the model.
 \item Verify that this partial $F$ statistic equals the square of the
 \texttt{hp} $t$ statistic from Exercise 8.
 \item Run \texttt{plot(full\_model)}. State what each of the four standard
 plots is designed to check.
 \item Identify the observations with the largest absolute standardized
 residual and largest Cook's distance. Should they be deleted automatically?
\end{enumerate}
\end{exercise}

\newpage
\section{Worked solutions}

\begin{answer}{1}
\begin{lstlisting}
head(cars)
dim(cars)
names(cars)
str(cars)
\end{lstlisting}
There are 32 observations and 11 variables. The three variables used here are
stored as numerical variables. Each row represents one car model.
\end{answer}

\begin{answer}{2}
\begin{lstlisting}
x <- cars$wt
y <- cars$mpg
x_bar <- mean(x)
y_bar <- mean(y)
s_x <- sd(x)
s_y <- sd(y)
c(x_bar=x_bar, y_bar=y_bar, s_x=s_x, s_y=s_y)
\end{lstlisting}
The results are
\[
 \bar x=3.21725,
 \qquad \bar y=20.090625,
 \qquad s_x=0.978457,
 \qquad s_y=6.026948.
\]
Next,
\begin{lstlisting}
Sxx <- sum((x - x_bar)^2)
Sxy <- sum((x - x_bar) * (y - y_bar))
b1 <- Sxy / Sxx
b0 <- y_bar - b1 * x_bar
c(Sxx=Sxx, Sxy=Sxy, b0=b0, b1=b1)
\end{lstlisting}
This gives
\[
 S_{xx}=29.678748,
 \qquad S_{xy}=-158.617225,
\]
\[
 \widehat\beta_0=37.285126,
 \qquad \widehat\beta_1=-5.344472.
\]
\end{answer}

\begin{answer}{3}
\begin{lstlisting}
simple_model <- lm(mpg ~ wt, data = cars)
coef(simple_model)
\end{lstlisting}
The output is
\begin{lstlisting}
(Intercept)          wt
  37.285126   -5.344472
\end{lstlisting}
Therefore
\[
 \widehat{\mathrm{mpg}}=37.2851-5.3445\,\mathrm{wt}.
\]
A 1000-pound increase in weight is associated with an estimated decrease of
about 5.34 mpg. This is an association, not automatically a causal effect.
The coefficients agree with the manual least-squares calculation.
\end{answer}

\begin{answer}{4}
\begin{lstlisting}
full_model <- lm(mpg ~ wt + hp, data = cars)
coef(full_model)
\end{lstlisting}
R returns
\begin{lstlisting}
(Intercept)          wt          hp
37.22727012 -3.87783074 -0.03177295
\end{lstlisting}
Thus
\[
 \widehat{\mathrm{mpg}}
 =37.2273-3.8778\,\mathrm{wt}-0.031773\,\mathrm{hp}.
\]
Holding horsepower fixed, an additional 1000 pounds is associated with about
3.88 fewer mpg. Holding weight fixed, one additional horsepower is associated
with about 0.0318 fewer mpg. The weight coefficient changes from $-5.3445$ to
$-3.8778$ because weight and horsepower contain overlapping information; a
multiple-regression coefficient isolates the contribution of one predictor
conditional on the other.
\end{answer}

\begin{answer}{5}
\begin{lstlisting}
y_hat <- fitted(full_model)
e <- resid(full_model)
max(abs(cars$mpg - (y_hat + e)))
sum(e)
sum(cars$wt * e)
sum(cars$hp * e)
\end{lstlisting}
The first result is exactly 0 in this run. The remaining results are
approximately
\[
 3.44\times10^{-15},\quad
 4.47\times10^{-15},\quad
 2.86\times10^{-13}.
\]
Mathematically, all three are zero. Computers store most decimal values with
finite binary precision, so tiny rounding errors remain. They should be read as
numerical zero, not as violations of the least-squares identities.
\end{answer}

\begin{answer}{6}
\begin{lstlisting}
X <- model.matrix(full_model)
Y <- matrix(cars$mpg, ncol = 1)
beta_matrix <- solve(t(X) %*% X, t(X) %*% Y)
beta_matrix
coef(full_model)
\end{lstlisting}
Both methods give
\[
 \widehat{\boldsymbol\beta}
 =\begin{pmatrix}
 37.22727012\\-3.87783074\\-0.03177295
 \end{pmatrix}.
\]
The check
\begin{lstlisting}
drop(t(X) %*% matrix(e, ncol = 1))
\end{lstlisting}
produces values on the order of $10^{-13}$ or smaller. Thus
$\mathbf X^{\mathsf T}\mathbf e=\mathbf0$: residuals are orthogonal to the
intercept column and both predictor columns. This condition is the matrix form
of the normal equations.
\end{answer}

\begin{answer}{7}
\begin{lstlisting}
n <- nrow(cars)
p <- 2
SSE <- sum(e^2)
MSE <- SSE / (n - p - 1)
SST <- sum((cars$mpg - mean(cars$mpg))^2)
R2 <- 1 - SSE / SST
c(SSE=SSE, MSE=MSE, residual_SE=sqrt(MSE), R2=R2)
\end{lstlisting}
The results are
\[
 \mathrm{SSE}=195.047755,
 \qquad
 \mathrm{MSE}=6.725785,
\]
\[
 \sqrt{\mathrm{MSE}}=2.593412,
 \qquad
 R^2=0.826786.
\]
The error degrees of freedom are $32-2-1=29$. The residual standard error and
$R^2$ agree with \texttt{summary(full\_model)}. The adjusted value is
$R^2_{\mathrm{adj}}=0.814840$.
\end{answer}

\begin{answer}{8}
\begin{lstlisting}
tab <- coef(summary(full_model))
t_hp <- tab["hp", "Estimate"] / tab["hp", "Std. Error"]
df_error <- df.residual(full_model)
p_hp <- 2 * pt(abs(t_hp), df=df_error, lower.tail=FALSE)
c(t_hp=t_hp, df=df_error, p_value=p_hp)
confint(full_model, "hp", level=0.95)
\end{lstlisting}
For horsepower,
\[
 \widehat\beta_{\mathrm{hp}}=-0.03177295,
 \qquad
 \mathrm{SE}(\widehat\beta_{\mathrm{hp}})=0.00902971.
\]
Therefore
\[
 t=\frac{-0.03177295}{0.00902971}=-3.518712,
 \qquad df=29,
\]
and the two-sided $p$-value is $0.001451$. The 95\% confidence interval is
\[
 (-0.050241,-0.013305).
\]
Because $p<0.05$, reject $H_0$. There is evidence that horsepower is associated
with mpg after controlling for weight. The confidence interval supports the
same conclusion because it excludes zero.
\end{answer}

\begin{answer}{9}
\begin{lstlisting}
new_car <- data.frame(wt=3, hp=110)
predict(full_model, newdata=new_car,
        interval="confidence", level=0.95)
predict(full_model, newdata=new_car,
        interval="prediction", level=0.95)
\end{lstlisting}
The predicted value is $22.09875$ mpg. The 95\% confidence interval for the
mean response is
\[
 (21.01879,23.17871),
\]
while the 95\% prediction interval for one new car is
\[
 (16.68580,27.51170).
\]
The prediction interval is wider because it includes both uncertainty in the
estimated mean and the natural car-to-car error $\varepsilon_{\mathrm{new}}$.
\end{answer}

\begin{answer}{10}
\begin{lstlisting}
reduced_model <- lm(mpg ~ wt, data=cars)
anova(reduced_model, full_model)
\end{lstlisting}
The comparison reports
\[
 F=12.381,
 \qquad p=0.001451.
\]
Adding horsepower significantly reduces SSE, so the full model is preferred by
this test. Also
\[
 (-3.518712)^2=12.381,
\]
confirming $F=t^2$ when one coefficient is added.

Run the standard plots with
\begin{lstlisting}
par(mfrow=c(2, 2))
plot(full_model)
par(mfrow=c(1, 1))
\end{lstlisting}
Their main purposes are:
\begin{enumerate}[leftmargin=*]
 \item \textbf{Residuals vs Fitted:} check linearity and look for changing
 spread.
 \item \textbf{Normal Q--Q:} check whether residuals are approximately normal.
 \item \textbf{Scale--Location:} check constant error variance.
 \item \textbf{Residuals vs Leverage:} find potentially influential cases.
\end{enumerate}
Observation 20 (Toyota Corolla) has the largest absolute standardized residual,
and observation 17 (Chrysler Imperial) has the largest Cook's distance. A flag
means ``investigate,'' not ``delete.'' Removal requires a data-quality or
scientific justification.
\end{answer}

\newpage
\section*{One-page R regression checklist}
\addcontentsline{toc}{section}{One-page R regression checklist}

\begin{lstlisting}
# 1. Inspect
head(d); str(d); summary(d)

# 2. Fit and read
fit <- lm(y ~ x1 + x2, data=d)
summary(fit)
coef(fit)
confint(fit)

# 3. Extract
y_hat <- fitted(fit)
e <- resid(fit)
X <- model.matrix(fit)

# 4. Verify the math
sum(e)                       # approximately 0
t(X) %*% e                  # approximately zero vector
SSE <- sum(e^2)
MSE <- SSE / df.residual(fit)
R2 <- 1 - SSE / sum((d$y - mean(d$y))^2)

# 5. Predict
new_case <- data.frame(x1=..., x2=...)
predict(fit, new_case, interval="confidence")
predict(fit, new_case, interval="prediction")

# 6. Compare and diagnose
anova(reduced_fit, full_fit)
par(mfrow=c(2, 2)); plot(fit)
\end{lstlisting}

\noindent\textbf{Final memory rule:}
\[
 \boxed{\text{formula} \longrightarrow \texttt{lm()}
 \longrightarrow \texttt{summary()}
 \longrightarrow \text{residual checks}}
\]

\section*{Common mistakes}
\begin{itemize}[leftmargin=*]
 \item Use \texttt{==} to test equality; use \texttt{<-} to assign a value.
 \item Use \texttt{\%*\%} for matrix multiplication, not \texttt{*}.
 \item Put new predictor values in a data frame with exactly the original
 variable names before calling \texttt{predict()}.
 \item In $n-p-1$, $p$ counts predictors but not the intercept.
 \item Do not interpret a multiple-regression slope without saying ``holding
 the other predictors fixed.''
 \item A small $p$-value does not measure practical importance or prove
 causation.
 \item Diagnostic flags do not automatically justify deleting observations.
\end{itemize}

\end{document}
